% ============================================================
% sections/related.tex  関連研究（初稿 2026-10-26）
% 方針：新規性を言い過ぎない。仮説を比べる推論そのものは新しくないと明記する。
%   機械学習側は短く（逆計画と Raileanu らは必ず触れる）。発達研究は類推にとどめる。
% ============================================================
\section{関連研究}\label{sec:related}

\subsection{シミュレーション説と自他の共有表現}\label{sec:related-st}
シミュレーション説は，他者の心的状態を，自分の心の働きを他者の立場で動かすことで推し量ると考える\cite{gordon1986folk,goldman2006simulating}．
神経科学では，自分の行為・感覚と他者のそれとに共通する表現（共有表現）と，両者を区別する仕組みの両方が社会的認知に必要だと整理されている\cite{decety2003shared}．
計算論の側では，運動制御のための内部モデルを他者の行為の理解や模倣に用いる枠組み\cite{wolpert2003unifying}や，ミラーニューロン系の働きを予測符号化として捉える考え方\cite{kilner2007predictive}が提案され，シミュレーションと明示的な心的推論を統合する議論もある\cite{keysers2007integrating}．
Gordon\cite{gordon2021simulation}は，予測符号化の観点から，自分の意思決定や計画の仕組みが他者の行動の仮説的な説明を試すためにも使われ，その際の信念や欲求は自他を区別しない符号化（agent-neutral coding）で表されると論じている．
重ね合わせ機構\cite{noguchi2022superposition}は，こうした考え方を，自他の区別を与えない予測学習の計算モデルとして具体化したものと位置づけられる．
本研究は，この機構の上で，自分の価値の経験がどこまで他者に当てはめられるかを調べる．

\subsection{投影と自己中心性バイアス}\label{sec:related-projection}
人は，自分の判断や好みを他者も共有していると過大に見積もる（偽の合意効果\cite{ross1977false}）．
自分の知識を他者に当てはめる傾向\cite{nickerson1999how}や，空腹などその時の身体の状態を他者に投影する傾向\cite{vanboven2003social}も報告されている．
他者の視点を取ることは，自分の視点を出発点にして調整する過程として説明されている\cite{epley2004perspective}．
Tamir と Mitchell\cite{tamir2010neural}は，他者の好みを推し量るときの内側前頭前野の活動が，自分の好みと他者についての判断の隔たりの大きさに応じて変わることを示し，人は自分を出発点にしてそこから調整すると結論した．
また，自分に似た他者について推し量るときには，自分についての思考に関わる脳部位が働く\cite{mitchell2006dissociable}．

シミュレーション説は，このような自己中心的な誤りを，自分の心の状態が他者についての判断に混ざることの帰結として説明する\cite{goldman2006simulating}．
一方 Saxe\cite{saxe2005against}は，他者の心の推し量りに見られる誤りを根拠に，シミュレーション説に反対する議論を行った．
% 【要確認】Saxe (2005) の議論の中身（どの型の誤りを挙げているか）は本文で確かめてから、必要なら一文足す。
本研究は，この論争に決着をつけるものではない．
本研究が示すのは，自他で共有された予測の仕組みが，当てはめと同時に，行動の仕方が自分と違う他者への投影の誤りに当たる予測の外れを生みうることである．

\subsection{行動から目標や好みを推す研究}\label{sec:related-inverse}
他者の行動からその目標や好みを推す問題は，認知科学と機械学習で広く研究されてきた．
逆計画\cite{baker2009action}は，他者がおおむね合理的に行動すると仮定し，候補となる目標ごとに予想される行動と観察を比べて目標を推す．
この枠組みは信念や知覚の推論にも広げられ\cite{baker2017rational}，人が他者の選択からその好みを逆向きに学ぶことも示されている\cite{jern2017people}．
機械学習では，観察された行動から報酬関数を推す逆強化学習\cite{ng2000algorithms}や，ニューラルネットワークに心の理論を学習させる研究\cite{rabinowitz2018machine}がある．

本研究の実験4（仮説の照合）も，候補ごとの予測を観察と比べる点で，これらと同じ構造を持つ．
したがって，候補を比べる推論そのものを新しい手法として主張はしない．
違いは，比べるための予測モデルを，他者の行動データや与えられた合理性の仮定からではなく，自分の価値の経験の予測学習から獲得し，共有モジュールを通して他者に当てはめている点にある．
また候補は研究者が用意した二択であり，好みを自ら見いだすことではなく，与えられた仮説を選ぶことにあたる．

\subsection{自分をもとに他者をモデル化する研究}\label{sec:related-self}
自分の価値の計算を他者に流用する仕組みは，人の脳でも調べられている．
Suzuki ら\cite{suzuki2012learning}は，他者の価値に基づく選択を予測する課題で，他者について自分の中で模擬した価値と，他者が実際に得た報酬との差を表す信号が，自分自身の価値づけの信号と重なる腹内側前頭前野に現れることを示した．
同時に，模擬から予測した他者の行動と実際の行動との差を表す信号が背内側前頭前野に現れ，この行動の予測誤差による補正は，他者の価値が自分と異なるときに重要になると示唆している．
本研究のモデルは，この区別でいえば，自分の価値づけを流用する経路だけを持ち，他者の行動の予測誤差によって他者の価値を補正する専用の経路を持たない．
本研究は，その状態で何が生まれ，何が生まれないかを示すものと位置づけられる．

計算モデルで最も近いのは，自分の方策を用いて他者の目標を推す Raileanu ら\cite{raileanu2018modeling}の研究である．
そこでは，自分の目標を入力の一つとする自分の方策を他者に当てはめ，他者の観察された行動がもっともらしくなるように目標の入力を直接最適化することで，他者の隠れた目標を推す．
本研究との違いは三点ある．
第一に，彼らの自分のモデルは，目標に依存する報酬を最大化する強化学習の枠組みで得た方策であり，目標が入力として明示されている．これに対し本研究の見込みは，報酬も目標の入力も用いない予測学習の中で，自分の行動の経験から生まれる．
第二に，本研究は自分の経験の量と自他の重みの共有を操作し，当てはめがそれらにどう依存するかを対照で示す．
第三に，本研究は当てはめの成功だけでなく，投影の誤りと，他者用のモジュールが好みを書き出さないことという限界も示す．

構成論的な研究では，Copete ら\cite{copete2016motor}が，視覚・触覚・関節角を予測する一つの予測器を自分の行為と他者の行為の観察の両方に用い，自分の運動の学習が進むほど他者の行為の結果を予測できるようになることを示した．
自分の行為の経験が同じ予測器を通して他者の予測を助けるという点で本研究と近い．
ただし Copete らは，観察者と行為者が同じ視点を共有すると仮定して自他の対応を与えており，他者を観察するときは，同じ予測器が自分の経験から欠けた感覚を補う．
% 確認：Copete ら (2016) の II-B 節（同じ視点の仮定）・VI 節（自他の対応は今後の課題）。研究室公開の PDF で確認（2026-10-26）。
これに対し重ね合わせ機構では，自他の対応は予測学習から生まれ\cite{noguchi2022superposition}，他者のための内部状態が自分の内部状態とは別に保たれる．
本研究ではその対応を土台として固定して用いるが，この構成には二つの利点がある．
一つは，他者の内部状態に自分の見込みが入り込むので，投影がその場で起き，その強さと代償を測れることである．
もう一つは，自他の重みの共有だけを外す対照を作れることである．
加えて，本研究は価値に従う見込みを扱い，その強さを経験の量で操作する．

重ね合わせ機構の系譜では，共有モジュールが視覚と運動から自分の位置の地図を自己組織化すること\cite{noguchi2017cognitive}，複数の感覚を一つの共有モジュールで扱うと地図の表現と地図の読み取りが生まれること\cite{noguchi2022multi}が示されている．
同じグループの Auchi ら\cite{auchi2022understanding}は，重ね合わせ機構に基づく視覚予測学習によって他者の意図を理解することを扱った．
% 【要確認】Auchi ら (2022) の「意図」が何を指すか（他者の行動の型か、目標か）は予稿の本文で確かめる。
%   確かめられたら、違いを一文で具体的に書き直す。
本研究が扱うのは，他者の意図や目標そのものではなく，価値（快・不快）に従う選択の見込みである．
また，その見込みの強さを自分の経験の量で操作し，共有を除く対照と，投影の誤りおよび他者用のモジュールの限界を示す点が異なる．
本研究は，この系譜の中で，共有モジュールを自他の価値の扱いへ広げる試みである．

\subsection{発達研究との対応}\label{sec:related-dev}
乳児は，他者が手を伸ばす行為の目標を，手が向かう場所ではなく物として選択的に符号化する\cite{woodward1998infants}．
さらに，3か月児のうち，物がくっつくミトンで自分で物をつかむ経験を先にした乳児だけが，その後に見た他者の行為者と目標の関係に注意を向けた\cite{sommerville2005action}．
自分の行為の経験が，他者の行為の理解を変えるのである．
一方，他者の好みの理解では，14か月児は，実験者が自分とは違う食べ物を好むことを表情で示しても，自分の好む食べ物を差し出す自己中心的な反応を示した\cite{repacholi1997early}．
% 同じ研究の18か月児の結果（相手の好みに沿って渡す）は再現性が議論されているため引用の中心にしない（執筆メモ §10、2026-10-26）。
%   内容は Boston College の講義資料に再掲された原著の要旨と結果の記述で確認した（一次の要旨の頁は読めなかった）。

本研究の環境と条件は，これらの課題と次のように対応づけられる．
物を色（赤と緑のランドマーク）に，置き場所の入れ替えを赤と緑の配置の入れ替えに，自分でつかむ経験を価値に従って動いた経験に対応させる．
また，自分の好みを他者に当てはめる14か月児の反応は，本研究の投影の誤りと，他者用のモジュールが他者固有の好みを表さないことに対応する．
ただし，これは計算モデルの振る舞いと発達の現象との類推であり，本研究は人間の発達について何かを示すものではない．
